\section{Related Work}
\label{sec:related}

\paragraph{Linear representations in LLMs.}
The linear representation hypothesis~\cite{park2023linear} posits that high-level concepts are encoded as linear directions in a model's activation space. Empirical support comes from diverse domains: \citet{marks2023geometry} show that truth and falsehood are linearly separable in LLM hidden states, \citet{burns2022discovering} recover latent knowledge via unsupervised linear probes, and \citet{li2023inference} demonstrate that linear interventions on ``truthfulness directions'' improve factual accuracy. \citet{nanda2023emergent} find that even game-playing models develop linear board-state representations. In safety, \citet{arditi2024refusal} show that refusal behavior across 13 models (up to 72B parameters) is mediated by a single linear direction, enabling both ablation and induction of refusal via weight orthogonalization. These findings suggest that many critical model properties are geometrically simple and, we hypothesize, therefore accessible to introspection.

\paragraph{Nonlinear and multi-dimensional features.}
Not all features are one-dimensional. \citet{engels2025not} provide a theoretical framework for \emph{irreducible multi-dimensional features}---representations that cannot be decomposed into independent one-dimensional components. They discover circular features (days of the week, months) in GPT-2 and Mistral~7B, showing these are causally used for modular arithmetic. Most relevant to our work, \citet{hildebrandt2025refusal} demonstrate that refusal behavior in LLMs is fundamentally nonlinear: UMAP and t-SNE achieve substantially better cluster separability than PCA across six models and three architecture families, with sub-clusters emerging within harmful-content embeddings that are invisible to linear methods. This tension between the linear view of \citet{arditi2024refusal} and the nonlinear view of \citet{hildebrandt2025refusal} motivates our hypothesis: binary refusal may be linear, but the fine-grained structure within refusal is not.

\paragraph{Representation engineering and steering.}
\citet{zou2023representation} introduce representation engineering as a top-down approach to reading and controlling high-level concepts (honesty, emotion, fairness) via linear directions in activation space. This framework implicitly assumes that features worth monitoring are linearly accessible. Our work tests this assumption: if some safety-relevant features are nonlinear, representation engineering may miss them, and models may be unable to self-report on them.

\paragraph{Model self-knowledge and introspection.}
A growing body of work examines whether LLMs can accurately report on their own internal states. \citet{tamoyan2025factual} find that models encode linear features predicting whether they will recall correct facts---a form of ``factual self-awareness'' that emerges rapidly during training and peaks in intermediate layers. \citet{yang2025when} show that retraction behavior (admitting mistakes) is driven by linearly probeable ``belief'' states that causally influence model outputs. These studies demonstrate that LLMs possess some introspective capability, but both focus exclusively on linearly encoded properties. We extend this line of inquiry by asking whether introspective access degrades for nonlinearly encoded features.

\paragraph{Probing methodology.}
Our approach builds on the probing paradigm reviewed by \citet{belinkov2021probing}, which highlights that the complexity of the probe matters: a sufficiently powerful probe can extract information that a simpler one cannot. We exploit this property by comparing linear (logistic regression) and nonlinear (MLP) probes on the same features, using the accuracy gap as a measure of feature nonlinearity. This follows the logic that if an MLP substantially outperforms logistic regression, the feature contains structure that is not linearly accessible~\cite{belinkov2021probing}.
